\section{Statement of results}\label{s:classification}
In this section we more formally introduce the Racah algebra~$\Re$ and state the main result of the paper which gives a classification of the finite-dimensional irreducible modules of the Racah algebra~$\Re$. This main result will be proved later in Section~\ref{s:proof}.

\begin{Definition}[{\cite[Definition 3.1]{SH:2017-1}}] \label{defn:Racah}
The {\it Racah algebra} $\Re$ is the unital associative $\F$-algebra defined by generators and relations in the following way. The generators are $A$, $B$, $C$, $D$. The relations state that
\begin{gather}\label{r:D}
[A,B]=[B,C]=[C,A]=2D
\end{gather}
and that each of
\begin{gather*}
[A,D]+AC-BA, \qquad [B,D]+BA-CB, \qquad[C,D]+CB-AC
\end{gather*}
is central in $\Re$.
\end{Definition}
It follows from the above definition that the element $A+B+C$ is also central in $\Re$. For notational convenience, we let
\begin{gather}
\alpha = [A,D]+AC-BA, \label{alpha} \\
\beta = [B,D]+BA-CB,\label{beta}\\
\gamma = [C,D]+CB-AC,\label{gamma}\\
\delta = A+B+C. \label{delta}
\end{gather}

\begin{Lemma}\label{lem:delta}\quad
\begin{enumerate}\itemsep=0pt
\item[$(i)$] The Racah algebra $\Re$ is generated by the elements $A$, $B$, $C$.
\item[$(ii)$] The Racah algebra $\Re$ is generated by the elements $A$, $B$, $\delta$.
\end{enumerate}
\end{Lemma}
\begin{proof}
(i) By~(\ref{r:D}) the element $D$ can be expressed in terms of $A$, $B$. Hence (i) follows from Definition~\ref{defn:Racah}.
(ii) By (\ref{delta}) the element $C$ can be expressed in terms of $A$, $B$, $\delta$. Hence (ii) follows from~(i).
\end{proof}

\begin{Lemma}\label{lem:presentationR}The $\F$-algebra $\Re$ has a presentation given by generators $A$, $B$, $\alpha$, $\beta$, $\delta$ and relations
\begin{gather}
A^2 B-2 A B A+B A^2-2AB-2BA=2A^2-2A\delta +2\alpha,\label{AAB}\\
AB^2-2BAB+B^2 A-2 AB-2 BA= 2B^2-2B\delta-2\beta,\label{ABB}\\
\alpha A = A\alpha, \qquad\beta A = A\beta, \qquad\delta A = A\delta,\notag\\
\alpha B = B\alpha,\qquad\beta B = B\beta,\qquad\delta B = B\delta,\qquad
 \alpha \delta = \delta \alpha,\qquad \beta \delta = \delta \beta. \notag
\end{gather}
\end{Lemma}
\begin{proof}We know from Lemma~\ref{lem:delta}(ii) that $A$, $B$, $\alpha$, $\beta$, $\delta$ generate~$\Re$. Observe that $C=\delta-A-B$ by \eqref{delta} and $D=\frac{1}{2}[A,B]$ by~\eqref{r:D}. The result can now be obtained by either using these two facts to eliminate~$C$, $D$ from the presentation of $\Re$ given in Definition~\ref{defn:Racah} or by using $D=\frac{1}{2}[A,B]$ to eliminate~$D$ from the presentation of $\Re$ given in \cite[Proposition~3.4]{SH:2017-1}.
\end{proof}

In the following proposition, we assert the existence of certain finite-dimensional $\Re$-modules and describe the actions of the generators of $\Re$ on these modules. A reader familiar with the theory of tridiagonal pairs will immediately recognize the form of the matrices representing~$A$ and~$B$ as precisely those given in Terwilliger's 2001 seminal work on tridiagonal pairs \cite[Theorem~3.2]{lp2001}.

\begin{Proposition}\label{prop:Rd}For any $a,b,c\in \F$ and any $d\in \N$ there exists a $(d+1)$-dimensional $\Re$-module $R_d(a,b,c)$ satisfying each of the following conditions:
\begin{enumerate}\itemsep=0pt
\item[$(i)$] There exists an $\F$-basis $\{v_i\}_{i=0}^d$ for $R_d(a,b,c)$ with respect to which the matrices represen\-ting $A$ and $B$ are
\begin{gather*}
\begin{pmatrix}
\theta_0 & & & &{\bf 0}
\\
1 &\theta_1
\\
&1 &\theta_2
 \\
& &\ddots &\ddots
 \\
{\bf 0} & & &1 &\theta_d
\end{pmatrix},
\qquad
\begin{pmatrix}
\theta_0^* &\varphi_1 & & &{\bf 0}
\\
 &\theta_1^* &\varphi_2
\\
 & &\theta_2^* &\ddots
 \\
 & & &\ddots &\varphi_d
 \\
{\bf 0} & & & &\theta_d^*
\end{pmatrix},
\end{gather*}
respectively, where
\begin{gather*}
\theta_i=\big(a+\textstyle\frac{d}{2}-i\big) \big(a+\textstyle\frac{d}{2}-i+1\big), \qquad 0\leq i\leq d,\\
\theta^*_i=\big(b+\textstyle\frac{d}{2}-i\big)\big(b+\textstyle\frac{d}{2}-i+1\big), \qquad 0\leq i\leq d,
\\
\varphi_i=i(i-d-1)\big(a+b+c+\textstyle\frac{d}{2}-i+2\big)\big(a+b-c+\textstyle\frac{d}{2}-i+1\big),
\qquad 1\leq i\leq d.
\end{gather*}

\item[$(ii)$] The elements $\alpha$, $\beta$, $\delta$ act on $R_d (a,b,c)$ as scalar multiplication by
\begin{gather*}
(c-b)(c+b+1)\big(a-\textstyle\frac{d}{2}\big)\big(a+\textstyle\frac{d}{2}+1\big),\\
(a-c)(a+c+1)\big(b-\textstyle\frac{d}{2}\big)\big(b+\textstyle\frac{d}{2}+1\big),\\
\textstyle\frac{d}{2} \big(\textstyle\frac{d}{2}+1\big)+a(a+1)+b(b+1)+c(c+1),
\end{gather*}
respectively.
\end{enumerate}
\end{Proposition}
\begin{proof} Using Lemma \ref{lem:presentationR}, this result can be verified through routine, though tedious, computations.
\end{proof}
